\chapter{Results and Analysis}
\label{sec:results}

This chapter is organized around the three thesis research questions. The archived broad and decoder-depth studies appear first because they explain how the final comparison set was chosen and which workflow variables were credible enough to carry forward. The repaired final held-out experiment then provides the main evidence for RQ1 through RQ3: what capability gaps remain on obfuscated binaries, which workflow variables materially affect performance, and whether completed analyses are properly backed by artifact evidence. The unrecovered final sweep is retained as capability-gap context, because it reveals operational failures that the repaired aggregate would otherwise hide.

\section{Development-Stage Archive Evidence}

The archived broad and decoder-depth figures are still useful, but they should be read as methodology-shaping artifacts rather than as final effect estimates. Their main value is that they show which workflow changes were worth carrying into the held-out evaluation and which ones exposed design problems before the blind final comparison was launched. The broad pilot sample set is summarized in Table~\ref{tab:broad-sweep-samples}, and the variant labels used across the archive and final studies are summarized in Table~\ref{tab:evaluation-variants}.

\subsection{Broad Coverage Pilot}

\begin{figure}[t]
\centering
\includegraphics[width=\textwidth]{experiment_results/broad_repaired_overall_ranking.png}
\caption{Archived broad-pilot ranking. This figure is used here as a descriptive development artifact: it shows which intervention families separated from the baseline strongly enough to justify follow-up study, but it is not treated as a thesis-grade estimate of general effectiveness.}
\label{fig:broad-repaired-ranking}
\end{figure}

Figure~\ref{fig:broad-repaired-ranking} supports two narrow conclusions. First, the pilot was sufficient to rule out some directions as uninteresting for the final comparison set. In particular, the final thesis sweep did not continue treating reduced tool access as a central variable, because that direction behaved more like an artificial handicap than an informative workflow comparison. Second, prompt-surface and worker-topology changes were large enough in the pilot to justify carrying them forward into deeper testing.

The broad pilot does \emph{not} support a stronger claim than that. It is best read as a screening step over the authored benchmark. It helps explain why the final sweep focused on worker subagents, worker prompt shape, response scope, and architectural emphasis, but it does not by itself establish that any of those variables generalize to harder held-out malware families.

\subsection{Focused Decoder-Depth Case Study}

\begin{figure}[t]
\centering
\includegraphics[width=\textwidth]{experiment_results/decoder_depth_combined_overall_ranking.png}
\caption{Archived decoder-depth ranking over the stripped and packed \texttt{config\_decoder\_test} variants. This case study is stronger than the broad pilot because it includes repeated runs, but it still remains a developmental case study rather than the final held-out evaluation.}
\label{fig:decoder-depth-ranking}
\end{figure}

\begin{figure}[t]
\centering
\includegraphics[width=\textwidth]{experiment_results/decoder_depth_combined_reliability.png}
\caption{Archived decoder-depth reliability and coverage. This figure is especially important because it separates weak analytical output from a validator-dominated failure mode.}
\label{fig:decoder-depth-reliability}
\end{figure}

The archived decoder-depth study supports a stronger but still limited interpretation. Compared with the broad pilot, it provides clearer evidence that workflow changes can matter differently on a deeper reverse-engineering task family than on easier overview-oriented tasks. Figure~\ref{fig:decoder-depth-ranking} shows that worker prompt shape and worker topology can change performance on packed and stripped decoder-oriented cases rather than merely changing report style. Figure~\ref{fig:decoder-depth-reliability} shows that not every additional layer helps. The validator-heavy pipeline behaves as a concrete negative case because its collapse is visible as a reliability and coverage failure rather than only a low score. The exact stripped and UPX-packed executables used in that case study are listed in Table~\ref{tab:depth-sweep-samples}, and the decoder-specific variant labels are defined in Table~\ref{tab:evaluation-variants}.

This case study therefore motivated two design choices for the held-out evaluation. The final sweep kept worker-topology and prompt-shape comparisons, because the decoder-depth figures suggested they could matter on richer tasks. It did not elevate the validator-heavy pipeline into the final comparison set, because the development evidence suggested that its failure mode was structural rather than merely unlucky on one sample. In other words, the archive studies do not answer the thesis research questions by themselves, but they do determine which variables were credible enough to carry forward into the held-out test that does.

\section{Final Held-Out Evaluation}

Unless otherwise noted, comparisons in this section refer to the rebuilt repaired experiment, which preserves the benchmark scope and configuration set while folding in successful recovery results for tasks that initially failed for infrastructure reasons. (The exact experiment identifier and artifact provenance are documented in Appendix~\ref{app:reproduction-commands}.) The held-out sample aliases and evaluator-side technique summaries are listed in Table~\ref{tab:final-heldout-samples}, and the configuration labels that appear in the figures are summarized in Table~\ref{tab:evaluation-variants}. Score-based claims in this section use the configured LLM-judge outputs as comparative proxy signals rather than as human labels or objective ground truth. The subsection order matches the research questions directly: RQ1 covers capability gaps on obfuscated binaries, RQ2 covers workflow-variable effects on performance, and RQ3 covers artifact-backed claims.

\subsection[RQ1: Capability Gaps on Obfuscated Binaries]{RQ1: What Are the Capability Gaps in an Agent's Ability to Perform Binary Analysis on Obfuscated Binaries?}

\begin{figure}[t]
\centering
\includegraphics[width=\textwidth]{experiment_results/final_eval_redundant_tool_usage_by_executable.png}
\caption{Per-executable semantic duplicate tool-call volume in the repaired final held-out experiment. The figure summarizes repeated tool activity within tasks, not proven waste. It is useful here because it highlights which held-out samples induce the heaviest repeated context reacquisition or tool reuse.}
\label{fig:final-redundant-tool-usage}
\end{figure}

RQ1 asks what the current system can and cannot do when the binaries are stripped, packed, encoded, or otherwise analysis-resistant. The held-out sample set is intentionally appropriate for that question: Table~\ref{tab:final-heldout-samples} shows that the evaluator-side reference cases combine stripping, packing, encoded configuration material, anti-analysis logic, staged behavior, and malware-family-inspired control-flow patterns. The repaired experiment suggests that the system can often produce usable static analyses on this kind of corpus, but it does not show universal or depth-complete success.

The most direct quantitative limit is task success. In the repaired final summary, the best-performing configurations, baseline and \texttt{artifact\_focused}, each reach a task-success rate of \texttt{0.867}; none reaches perfect completion across the benchmark. The single-generalist, blank worker-role, and Ghidra-focused variants each drop to \texttt{0.733}. Those numbers matter for RQ1 because the question is not only whether the system can produce a good report in the best case, but whether it can repeatedly stay operational long enough to finish analysis on obfuscation-heavy binaries.

One-shot robustness is a second major capability limit. The unrecovered held-out sweep exposed context-bloat and request-budget failures that were not dominant in the smaller development studies, especially on Ghidra-heavy worker assignments. That result matters because it shows a gap between ``can sometimes recover a useful answer'' and ``can reliably sustain analysis on harder obfuscation-heavy binaries.'' The repaired experiment is therefore the right artifact for comparing completed outputs, but the unrecovered sweep remains part of the answer to RQ1 because it exposes where the system stops being operational before analysis is complete.

Figure~\ref{fig:final-redundant-tool-usage} adds a third capability limit: repeated context reacquisition is concentrated on a subset of the held-out binaries rather than distributed evenly across the corpus. In total semantic duplicate volume, \texttt{sample5.exe} is the largest hotspot with \texttt{501} duplicate calls across the seven configurations, followed by \texttt{sample2.exe} with \texttt{364}. In the narrower same-source or likely-wasteful view discussed later in this chapter, \texttt{sample2.exe} and \texttt{sample3.exe} become the main hotspots. These counts do not prove that every repeated call was unnecessary, but they do show that some obfuscation-heavy or artifact-heavy executables are harder for the workflow to traverse cleanly than others.

The qualitative gap is therefore narrower than ``agents cannot analyze obfuscated binaries'' but wider than ``agents can reliably handle them.'' The repaired held-out comparisons show that the workflow can often recover broad behavior, configuration, persistence, or anti-analysis signals even on stripped and packed samples. At the same time, later-stage recovery, one-shot completion, and efficient handling of dense tool context remain weak points. RQ1 is therefore answered in a bounded way: the system is capable of useful static analysis on obfuscated binaries, but its main gaps remain reliable completion on harder cases, deeper stage recovery, and stable context management once tool outputs become large or repetitive.

\subsection[RQ2: Workflow Variables and Performance]{RQ2: What Workflow Variables Impact Performance on Analysis Capabilities?}

Throughout this subsection, performance means LLM-judge-based overall score together with task success and per-executable score deltas relative to the maintained baseline. Here the comparison point remains the inferred baseline carried through the experimental program: a tool-complete staged configuration with preparatory context gathering, deterministic presweeps, a balanced specialist worker mix, default prompt surfaces, and final reporting.

\begin{figure}[t]
\centering
\includegraphics[width=\textwidth]{experiment_results/final_eval_overall_ranking.png}
\caption{Repaired final held-out ranking across all seven configurations. This figure uses the rebuilt repaired experiment so that all configurations have full coverage. The highest descriptive mean score belongs to \texttt{worker\_prompt\_shape:artifact\_focused}, but none of the overall shifts reached the thesis confidence target of $p \le 0.05$.}
\label{fig:final-overall-ranking}
\end{figure}

\begin{figure}[t]
\centering
\includegraphics[width=\textwidth]{experiment_results/final_eval_effect_size_summary.png}
\caption{Effect-size summary for the repaired final held-out comparison relative to baseline. This figure matters because the repaired experiment supports descriptive ranking, but the inferential question is whether any overall shift is large and stable enough to treat as more than a sample-limited difference.}
\label{fig:final-effect-size-summary}
\end{figure}

\begin{figure}[t]
\centering
\includegraphics[width=\textwidth]{experiment_results/final_eval_cross_executable_heatmap.png}
\caption{Cross-executable score-delta heatmap for the repaired final held-out experiment. Positive cells indicate improvement over the baseline on that executable and negative cells indicate degradation. This view matters because every non-baseline configuration in the repaired final run is marked as mixed and masked by average in the executable-consistency summary.}
\label{fig:final-cross-executable-heatmap}
\end{figure}

Figures~\ref{fig:final-overall-ranking} and~\ref{fig:final-effect-size-summary} answer RQ2 at the aggregate level. Relative to baseline, the repaired final comparison does \emph{not} show a statistically significant overall winner at the thesis confidence target. The best descriptive overall mean score comes from \texttt{worker\_prompt\_shape:artifact\_focused} at \texttt{82.382}, compared with the baseline at \texttt{78.096}, but the replicate-level permutation test yields \texttt{p = 0.5}. Four of the six non-baseline variants underperform the baseline descriptively on overall mean score, and none of those degradations is statistically significant either. The aggregate answer to RQ2 is therefore mixed rather than triumphant: workflow variables do matter, but the repaired final run does not identify a statistically supported universal replacement for the maintained baseline.

Figure~\ref{fig:final-cross-executable-heatmap} shows why global averages are not enough. All six non-baseline configurations are heterogeneous across the five executables. In the executable-consistency summary, every non-baseline row is labeled \texttt{mixed}, and every one is flagged as \texttt{masked\_by\_average=true}. That is not a rhetorical caveat; it is a direct property of the repaired final dataset. The paged per-executable score panels generated from the same repaired experiment retain the repetition-level variance detail and are included in Appendix~\ref{app:final-per-executable-panels}, but the heatmap is the more legible thesis-scale summary and is therefore the figure kept in the main text.

The concrete pattern is easy to summarize. Baseline remains strongest on \texttt{sample2.exe} and \texttt{sample5.exe}. \texttt{worker\_prompt\_shape:empty} is the top descriptive performer only on \texttt{sample1.exe}, and even there the gain over baseline is modest at \texttt{+1.19}. \texttt{worker\_prompt\_shape:artifact\_focused} is strongest on \texttt{sample3.exe}, where it improves mean score by \texttt{+21.427} relative to baseline. \texttt{architecture\_preset:model\_augmented} ties for the best score on \texttt{sample4.exe} and improves by \texttt{+11.903} there. The same configuration is also the worst performer on \texttt{sample1.exe}, with a delta of \texttt{-17.856}. These are sample-dependent reversals, not minor fluctuations around the same answer. Table~\ref{tab:final-heldout-samples} is useful alongside this subsection because the alias names alone do not tell the reader which kinds of behaviors or analysis challenges each sample was meant to stress.

Task-success and efficiency summaries reinforce the same conclusion. Baseline and \texttt{artifact\_focused} each reach a task-success rate of \texttt{0.867}, while the single-generalist, blank worker-role, and Ghidra-focused variants fall to \texttt{0.733}. Mean task wall-clock duration ranges from \texttt{1324.595} seconds for \texttt{single\_generalist} to \texttt{1716.168} seconds for \texttt{artifact\_focused}. These movements show that workflow variables affect more than report style: they change completion rate, runtime, and sample-specific behavior. But the effects do not move in a single direction. The evidence-first worker-prompt overlay is the strongest descriptive candidate in the current design space, while the remaining alternatives either appear neutral-to-negative overall or depend heavily on which executable is being analyzed.

RQ2 is therefore answered in a constrained but concrete way. Workflow variables do impact performance on analysis capabilities, but their effects are conditional. Prompt shape, worker mix, and architectural emphasis can each help or hurt depending on the executable profile, and the current repaired final dataset supports executable-specific findings more confidently than global ones. The next interpretive step is plausible but not fully proven: the executables that benefit from \texttt{artifact\_focused} or \texttt{model\_augmented} may be the ones where deeper control-flow reconstruction or evidence organization matters more than broad default coverage. With only five executables, however, that remains a hypothesis for future stratified evaluation rather than a settled causal explanation.

\subsection[RQ3: Artifact-Backed Claims]{RQ3: Are Agents Able to Back Claims with Proper Artifact Evidence?}

\begin{table}[t]
\centering
\scriptsize
\caption{Mean LLM-judge rubric-dimension scores most directly related to artifact-backed analysis in the repaired final held-out experiment. Values are mean rubric scores across three repetitions per configuration on the 1--5 judge scale.}
\label{tab:final-evidence-dimensions}
\begin{tabular}{|p{1.95in}|c|c|c|c|}
\hline
Configuration & Evidence grounding & Specificity & False-claim control & Task alignment \\
\hline
baseline & 3.867 & 3.933 & 3.600 & 4.267 \\
artifact\_focused prompt & 4.267 & 4.267 & 3.733 & 4.400 \\
model\_augmented architecture & 3.867 & 4.067 & 3.267 & 4.200 \\
ghidra\_focused architecture & 3.800 & 3.800 & 3.400 & 4.200 \\
single\_generalist workers & 3.733 & 3.800 & 3.267 & 3.933 \\
empty worker prompt & 3.667 & 3.933 & 3.067 & 4.067 \\
brief response scope & 3.600 & 3.867 & 3.400 & 4.067 \\
\hline
\end{tabular}
\end{table}

RQ3 asks whether the system is producing grounded reverse-engineering claims rather than only plausible summaries. Table~\ref{tab:final-evidence-dimensions}, derived from the repaired final held-out \texttt{dimension\_summary.csv}, gives the most direct aggregate proxy evidence under the LLM-judge rubric. The baseline itself is not evidence-free: it reaches \texttt{3.867} on evidence grounding, \texttt{3.933} on specificity, and \texttt{3.600} on false-claim control. That means completed baseline runs are often rated as at least moderately grounded by the configured judge. However, evidence quality is still configuration-sensitive. The evidence-focused prompt overlay is the strongest configuration on every dimension most relevant to RQ3, reaching \texttt{4.267} on evidence grounding, \texttt{4.267} on specificity, \texttt{3.733} on false-claim control, and \texttt{4.400} on task alignment. By contrast, the blank worker-role variant and the brief response-scope variant both reduce evidence-grounding quality relative to baseline.

The repaired \texttt{sample3.exe} runs provide a concrete illustration. In the baseline run set, the three repaired LLM-judge outputs score evidence grounding at \texttt{3}, \texttt{4}, and \texttt{4} out of \texttt{5}, with false-claim control at \texttt{2}, \texttt{3}, and \texttt{3}. One baseline replicate is summarized by the configured LLM judge as a ``decent overview'' that still ``lacked concrete evidence linking many technical claims to specific binary elements.'' In the matched \texttt{artifact\_focused} run set on the same sample, the three repaired LLM-judge outputs score evidence grounding at \texttt{5}, \texttt{5}, and \texttt{5}, with false-claim control at \texttt{5}, \texttt{5}, and \texttt{4}. The strongest repaired replicate is explicitly described by the configured LLM judge as covering the expected anti-analysis, persistence, and C2 behaviors with ``specific artifacts'' and with no unsupported claims recorded.

The answer to RQ3 is therefore positive but qualified under the configured LLM-judge rubric. Some completed runs are rated as artifact-backed, and the final held-out rubric scores show that some workflow choices improve that behavior. At the same time, evidence quality is not automatic. The baseline is only moderately grounded on average, some configurations erode false-claim control, and the strongest evidence-backed behavior appears on selected samples rather than as a guaranteed property of every completed run.

\subsection{Robustness Caveat from the Original Held-Out Sweep}
\label{sec:final-failure-analysis}

The repaired aggregate is the right artifact for comparing final scores, but the original unrecovered held-out sweep still matters as a methodological caveat. It exposed a real first-pass robustness problem concentrated on a subset of the harder samples, especially when Ghidra-heavy worker assignments accumulated too much shared context before completion. That matters because some one-shot failures were infrastructural rather than purely analytical, so the unrecovered sweep alone would understate the workflow's usable performance on the held-out benchmark.

That caveat should not dominate the results narrative, however. The main thesis question in this chapter is what can be learned from the completed final comparisons once recovery restores coverage. The unrecovered sweep is therefore best treated as context for interpreting those repaired comparisons, not as the central empirical result. The detailed outcome counts, per-executable failure figures, and assignment-failure breakdown are moved to Appendix~\ref{app:heldout-failure-accounting}, and Chapter~\ref{sec:limitations} discusses what they do and do not imply about capability gaps and external validity.

\section{Discussion and Claim Boundaries}

\subsection{Direct Answers to the Research Questions}

\paragraph{RQ1.} The system can often recover useful static findings on the held-out obfuscated binaries, but not with uniform reliability or depth. The repaired final comparison reaches at most \texttt{0.867} task success, while the unrecovered sweep exposed context-bloat and request-budget failures on harder cases. The best-supported answer is therefore that capability exists, but the main gaps remain one-shot completion, deeper staged recovery, and stable handling of large or repetitive tool context.

\paragraph{RQ2.} Workflow variables do affect performance, but their effects are conditional rather than universal. Relative to the maintained baseline, no tested configuration yields a statistically significant overall score improvement in the repaired final comparison. The strongest descriptive candidate is \texttt{worker\_prompt\_shape:artifact\_focused}, but the executable-level heatmap shows that every non-baseline variant is mixed across the five held-out samples.

\paragraph{RQ3.} The system can produce analyses that the configured LLM judge sometimes rates as artifact-backed, but not uniformly. The final held-out rubric summaries show moderately grounded baseline behavior and stronger evidence-grounding, specificity, and false-claim control under \texttt{artifact\_focused}. The repaired \texttt{sample3.exe} runs provide the clearest local example, but the same summaries also show that some configurations still weaken grounding quality or false-claim control.

\subsection{Plausible Interpretations}

The results also suggest, but do not prove, several interpretations. The first is that evidence-first prompt shaping may help most on executables where the bottleneck is organizing and grounding many competing artifacts rather than simply producing a concise report. The second is that architecture-level changes such as \texttt{model\_augmented} may be acting as conditional strategies rather than universal improvements: they help on some executables and clearly hurt on others. The third is that the held-out benchmark is stressing two capabilities at once, namely reverse-engineering quality and orchestration robustness. The failure analysis makes it plausible that some of the apparent performance difference between configurations is mediated by how well they keep worker contexts tractable, not only by how well they reason after the context is assembled. The efficiency summaries suggest a related operational hypothesis: some repeated work is likely caused by missing stage-to-stage reuse or state-loop controls rather than by analytical necessity alone.

\subsection{Open Uncertainties and Measurement Gaps}

Several uncertainties remain open. The final comparison still covers only five held-out executables and three repaired repetitions per configuration, so the executable-level patterns should be treated as strong case evidence rather than as a stable estimate of population-level behavior. The significance tables reinforce that caution: none of the overall score shifts reached the thesis significance threshold, and no multiple-comparison correction is applied in the exploratory task-level tables. In addition, the repaired experiment necessarily combines original runs with targeted recovery runs. That is appropriate for obtaining full benchmark coverage, but it means the final aggregate should be read as the best-available completed comparison under the evaluated workflow rather than as a pristine one-shot sweep.

The tool-usage summaries are also useful but limited. They identify repeated-call hotspots and they now distinguish same-source repeats from cross-source overlap within a task. However, they still do not attach every tool call to a specific work item inside a reused worker slot. That means the current instrumentation can support statements such as ``this task and source produced many repeated calls'' or ``same-source repetition exists here,'' but it cannot yet prove that each repeated call was truly unnecessary. The same repaired final study also uses a single \emph{default analysis} task per held-out executable and keeps the \texttt{repo\_default} model profile fixed across all seven configurations. As a result, the current efficiency analysis cannot cleanly separate task-type effects from executable effects or attribute redundancy differences to model-family changes. Those measurement limits are discussed more fully in Chapter~\ref{sec:limitations}.

Taken together, the repaired final results support three evidence-backed conclusions. First, the current workflow can often analyze obfuscation-heavy binaries usefully, but its capability gaps remain visible in completion failures, repeated context reacquisition, and incomplete later-stage recovery. Second, workflow-variable effects are real but conditional, so average-only score reporting would be misleading. Third, the workflow can produce analyses that the configured LLM judge rates as artifact-backed, but evidence quality is configuration-sensitive and can degrade under weaker prompt or role settings. These findings frame agentic reverse engineering as a capability-gap, workflow-selection, and evidence-grounding challenge rather than a prompt-tuning problem with a single optimal answer.
